\section{Benchmark Saturation}
\label{sec:benchmark-saturation}

The benchmark Humanity's Last Exam is introduced on the observation that the headroom MMLU relies on has closed: language models have now exceeded 90\% accuracy on popular benchmarks such as MMLU, which limits how much such a benchmark can still measure about state-of-the-art capability \cite{phanHumanitysLastExam2025}. Once most models under test score near the ceiling, the benchmark no longer separates them by capability, and the leaderboard collapses toward a narrow band of scores that no longer reflects differences in what the models can do.

Agent benchmarks report the same closing of headroom: top-performing agents have reached scores exceeding 90\% on existing benchmarks, which are becoming saturated and losing their ability to differentiate state-of-the-art agents \cite{guoSelfEvolvingBenchmarksSynthesizing2026}. Saturation is therefore a general property of predetermined benchmarks as models improve.

Benchmarks presuppose headroom: they separate models by capability only if the tasks are not already solved by every model under test. Headroom can close through mechanisms other than accumulated capability gains. Evaluation data can leak into training corpora, and fine-tuning on it inflates a model's score without a matching gain in capability \cite{akhtarWhenAIBenchmarks2026}. Benchmarks therefore need constant revision, proposing tasks difficult enough to preserve headroom \cite{akhtarWhenAIBenchmarks2026}.

Saturation is one way a predetermined evaluation stops discriminating between models. A second way is for the evaluation to be poorly matched to the domain it is applied to. One example is Singhal et al.'s work on medical question answering, where existing benchmarks are limited to classification accuracy or automated natural language generation metrics such as BLEU, and do not support the detailed analysis that real-world clinical use requires \cite{singhalLargeLanguageModels2023}. In place of such generic metrics, evaluation of clinical answers is carried out along axes specific to the domain, including factuality, comprehension, reasoning, possible harm, and bias \cite{singhalLargeLanguageModels2023}. A generic accuracy score or a generic text-generation metric does not capture what makes a clinical answer acceptable, so the metric is redefined around the domain instead of reused from a general-purpose benchmark.

A domain-tailored metric such as this one is still a single predetermined metric, applied to an output that is checked against a predetermined notion of what counts as acceptable. This is sufficient as long as the task itself has a comparison point of that kind. It stops being sufficient once the task is open-ended and has no predetermined answer to compare the output against, so that no single metric, however well matched to the domain, is left with anything to compare it to.
